% Generated draft: requires mathematical and source review.
\documentclass[11pt]{article}
\usepackage[margin=25mm]{geometry}
\usepackage{fontspec}
\setmainfont{Latin Modern Roman}
\usepackage{amsmath,amssymb,mathtools,unicode-math}
\setmathfont{Latin Modern Math}
\usepackage{enumitem,needspace,xurl,hyperref}
\setlength{\parindent}{0pt}
\setlength{\parskip}{5pt}
\setlength{\emergencystretch}{3em}
\newcommand{\N}{\mathbb{N}}
\newcommand{\Z}{\mathbb{Z}}
\newcommand{\Q}{\mathbb{Q}}
\newcommand{\R}{\mathbb{R}}
\newcommand{\C}{\mathbb{C}}
\newcommand{\F}{\mathbb{F}}
\newcommand{\A}{\mathbb{A}}
\newcommand{\PP}{\mathbb{P}}
\newcommand{\E}{\mathbb{E}}
\newcommand{\eps}{\varepsilon}
\newcommand{\dd}{\,\mathrm{d}}
\begin{document}
\section*{UnsolvedMath ID: 20000867}
\textbf{Source problem code:} catalogue entry
\subsection{Definitions and mathematical statement}
This source is Tai Sing Lee's Chapter J discussion, ``Resonance Prediction and Priors'', not a formal problem. It says V1--V2 response timing was hard to resolve, describes top-down predictions as more like a prior, and reports that V2 responses were not attenuated by attending elsewhere. A question proposes target removal after $100$ milliseconds; the speaker says this was not tried. The question about experiments with unknown target location has no recorded answer. The speaker attributes the longer pop-out latency, compared with bars, to three-dimensional surface inference. These are transcript statements, not a specified dynamical or causal model.

The aid proposes a conditional experimental-design certificate for that unperformed offset experiment. Let pre-offset history $H\in\{1,0\}$ distinguish an oddball-target condition from a control, and replace the target by a uniform distractor at $t_0=100\,\mathrm{ms}$ so the entire post-offset display is exactly matched. For a V1 or V2 response define its history contrast $\Delta(t)=\E[R(t)\mid H=1]-\E[R(t)\mid H=0]$. Declare a stimulus-memory horizon $M$ and timing uncertainty $\epsilon_t$, and evaluate a late window whose earliest time, reduced by $\epsilon_t$, exceeds $t_0+M$.

Suppose the estimators have simultaneous error radii $r(t)$ with probability at least $1-\alpha$, and the transition-artifact contribution to the contrast is certified bounded by $B(t)$. Writing $\Delta=\Delta_{\mathrm{res}}+\Delta_{\mathrm{art}}$, the aid's certificate is
\[
 |\Delta_{\mathrm{res}}(t)|\geq
 \max\{0,|\widehat\Delta(t)|-r(t)-B(t)\}
\]
on that simultaneous event. Thus $|\widehat\Delta(t)|>r(t)+B(t)$ rejects the declared bounded-memory null predicting no late residual history contrast. This requires matched displays, validated artifact bounds and an independently stated memory/timing margin. It does not prove that a residual is specifically a prediction, prior or feedback signal.

For a randomized and validated V2-to-V1 terminal intervention $Z\in\{0,1\}$, a factorial history/intervention design targets
\[
 \Psi=\bigl(\E[R\mid H=1,Z=1]-\E[R\mid H=0,Z=1]\bigr)
      -\bigl(\E[R\mid H=1,Z=0]-\E[R\mid H=0,Z=0]\bigr).
\]
With the stated randomization and intervention-validity assumptions, this difference-in-differences identifies the causal effect of that intervention on the history contrast. It does not identify the contrast's representational meaning or an entire circuit causal architecture; a null effect does not establish absence of feedback.

The summary adds three distinct limits of interpretation. First, onset estimates $\widehat T_1,\widehat T_2$ with certified absolute errors $e_1,e_2$ establish an order only if their gap exceeds $e_1+e_2$; even certified temporal order does not determine a causal arrow. Second, posterior-like response scores do not uniquely separate likelihood/prediction and prior contributions without an independently specified observation model. Third, a measured extra latency for three-dimensional stimuli supports a positive 3D-specific residual only after subtracting both timing errors and a bound on non-3D nuisance effects; a positive residual lower bound has the form $\widehat L_{3D}-\widehat L_{\mathrm{control}}-e_{3D}-e_{\mathrm{control}}-B_{\mathrm{nuisance}}$. No values or certified bounds are supplied. The proposed experiment and lemmas do not convert this contextual transcript into an already completed empirical test.
\subsection{Short English statement}
The transcript suggests an offset experiment but reports that it was not performed. The aid formalizes conditions under which a late history contrast would reject a bounded-memory explanation, and what a randomized feedback intervention could identify. Timing, prediction/prior labels and 3D-latency attribution require separate evidence.
\subsection{Sources}
\begin{itemize}
\item Local UnsolvedMath registry, permanent ID $20000867$, source code \texttt{catalogue entry}, complete supplied transcript and research summary; canonical local record \texttt{aim-biology-notes.json}, \texttt{notes[31]}.
\item American Institute of Mathematics, \emph{Inference and Prediction in Neocortical Circuits}, Chapter J, ``Resonance Prediction and Priors'', Tai Sing Lee: \url{https://aimath.org/WWN/brain/brain.pdf}. The supplied verification records the tenth physical PDF page, printed footer 11, as its checked source. No new PDF retrieval is claimed.
\item The aid identifies the discussion with high confidence as Lee et al.'s 2002 macaque shape-from-shading pop-out study, but supplies no full title, coauthor list or bibliographic identifier. This identification remains attributed rather than independently verified.
\item Research attempt 1: context only; problem not a problem at the run, registry status open and classification \texttt{OPEN-TRIAGE}; supplied rigor label ``proved'', novelty confidence ``low''. The machine-generated aid is not peer reviewed and requires independent verification before citation. No empirical dataset or authoritative citation for the proposed certificate is supplied.
\end{itemize}
\end{document}
